\documentclass[10pt,a4paper]{amsart}
\usepackage[margin=19mm]{geometry}
\usepackage{amsmath,amssymb,mathtools}
\usepackage[scr=rsfs,cal=euler]{mathalfa}
\usepackage{xfrac}
\usepackage[unicode,hidelinks,bookmarks=false]{hyperref}
\newcommand{\ie}{i.e.\ }
\newcommand{\cf}{cf.\ }
\newcommand{\resp}{respectively\ }
\newcommand{\Subsection}{\subsection}
\providecommand{\nobreakdash}{\nobreak\hbox{-}}
\setlength{\emergencystretch}{2em}
\multlinegap0pt
\usepackage{amssymb}
\usepackage{mathtools}
\newtheorem{theorem}{Theorem}
\newtheorem{lemma}{Lemma}
\title{A primality test for $Kp^n+1$ numbers and a generalization of Safe Primes and Sophie Germain Primes}
\author{A. Ramzy}
\date{Source: arXiv:2207.12407v1 (CC BY 4.0)}
\begin{document}
\maketitle

\noindent\emph{Source and licence.} A primality test for $Kp^n+1$ numbers and a generalization of Safe Primes and Sophie Germain Primes. Version arXiv:2207.12407v1 (CC BY 4.0), distributed by arXiv under the Creative Commons Attribution 4.0 International licence, http://creativecommons.org/licenses/by/4.0/, as recorded in the arXiv licence metadata for this exact version and shown on the abstract page https://arxiv.org/abs/2207.12407v1. Full licence notice: \texttt{licenses/arxiv-2207.12407-CC-BY-4.0.txt}.\par\smallskip

\noindent\textit{Editorial scope.} Counted unit: the article's theorem thm34 (source0.tex lines 275--282). The article's complete original proof of it is delivered with it (source0.tex lines 284--313, one \texttt{proof} environment of 30 lines). No mathematics is written, repaired or re-derived: the statement and the proof are the article's own lines, in the article's own notation, and no retained line is substituted or re-flowed.  Retained uncounted as the source-local context the proof needs: lines 223--230; lines 260--262.  Nothing was omitted from any retained span, so the package carries no omission record.  No retained line contains a citation, so the wrapper carries no bibliography and no bibliography entry.  No retained line contains a figure environment, an image-inclusion command or drawing code, so no image is delivered and the package declares no asset dependency.  Editorial formatting only, in addition: an AMS \texttt{amsart} preamble that restates the article's own declarations and macros used by the retained lines, namely the environments \texttt{theorem}, \texttt{lemma} on separate counters, together with the standard packages those lines need.  The retained environments print in this document as Theorem 1, Lemma 1 in order of appearance, whereas the article prints the same items with the numbering of its own full document.  The complete native source of the article as distributed in the arXiv e-print for this exact version is bundled unchanged as \texttt{sources/128708/source0.tex}.
\begin{theorem}\label{thm31}
	Let $N=Kp^n+1 \neq p^{3(n-j)}+1$,  where $p$ is prime and $p^j \ge p^{n-j} \ge (N-1)^{1/3}$. If there exists an integer $a>1$ such that: 
	\begin{enumerate}
		\item[(i)] $a^{Kp^{n-j-1}} \equiv L \neq 1 \pmod N$, and
		\item[(ii)] $L^{p^{j+1}} \equiv 1 \pmod N$.
	\end{enumerate}
	Then $N$ is prime.
\end{theorem}

\begin{lemma}\label{lem33}
	If $X > 11+4\sqrt{7}$, $a \le \frac{X}{2}$, $a \neq 1$, $a \neq \sqrt{X+1}$ and $\frac{(X-a)(a)+1}{X}=M \in \mathbb{Z}$. then $M > \sqrt{X}$.  
\end{lemma}

\begin{theorem}\label{thm34}
	Let $N=Kp^n+1 \neq p^{3(n-j)}+1$, $N \neq (\sqrt{p^{n-j}+1}-1)p^{3(n-j)}+1$ and $N \neq H^3p^{3(n-j)}+1$, where $p$ is prime, $p^j \ge p^{2(n-j)} $, $p^{n-j} \ge (N-1)^{2/7}$ and $p^{n-j} > 11+4\sqrt{7}$. If there exists an integer $a>1$ such that:
	\begin{enumerate}
		\item[(i)] $a^{Kp^{n-j-1}} \equiv L \neq 1 \pmod N$, and
		\item[(ii)] $L^{p^{j+1}} \equiv 1 \pmod N$.
	\end{enumerate}
	Then $N$ is prime.
\end{theorem}

\begin{proof}
	As in \ref{thm31}, we can deduce that there must exist a prime divisor $q$ of $N$ such that $p^{n-j} \mid q-1$, therefore, we have two cases:
	\begin{enumerate}
		\item[(i)] $q=N$ and $N$ is prime, or
		\item[(ii)] $q < N$ and $N$ is composite.
	\end{enumerate}
	Now, we will try to exclude (ii) by contradiction. In(ii) we have $q \mid N$, hence, $\frac{N}{q}=M \in \mathbb{Z}$, but since $q=ap^{n-j}+1$ then we must have $M=bp^{n-j}+1$ as well, now, if we multiply $q$ by $M$ we will get $N=Kp^n+1=abp^{2(n-j)}+(a+b)p^{n-j}+1 \Longrightarrow Kp^{j}=abp^{n-j}+a+b$. Therefore we must have $p^{n-j} \mid a+b$ and $p^{n-j} \mid (ab+\frac{a+b}{p^{n-j}})$ which leads us to the following two cases (let $p^{n-j}=X$ for abbreviation):
	\begin{enumerate}
		\item[(1)] $a,b <X \Longrightarrow a+b=X \Longrightarrow X \mid (ab+1) \Longrightarrow \frac{ab+1}{X}=\frac{N-1}{X^3}$
		\item[(2)] $a=AX+c$, $b=X-c \Longrightarrow  X \mid (AX^2 - AcX + cX -c^2 + A + 1) \Longrightarrow AX - Ac + c + \frac{A + 1 - c^2}{X}=\frac{N-1}{X^3}$
	\end{enumerate}
	In (1) we have $\frac{ab+1}{X} \in \mathbb{Z}$, and by assuming that $a \le b \Longrightarrow \frac{(X-a)(a)+1}{X} \in \mathbb{Z}$, therefore and as in \ref{lem33} (since $X > 11+4\sqrt{7}$) we can deduce that $a=1 \Longrightarrow \frac{N-1}{X^3}=1$ or $a=\sqrt{X+1} \Longrightarrow \frac{N-1}{X^3}=\sqrt{X+1}-1$ or $a > \sqrt{X}+2 \Longrightarrow \frac{(X-a)(a)+1}{X} = \frac{N-1}{X^3} > \sqrt{X}$ and the latter is impossible since $X \ge (N-1)^{2/7}$, hence, we have excluded all the cases of (1) (since we assumed that $N \neq p^{3(n-j)}+1$ and $N \neq (\sqrt{p^{n-j}+1}-1)p^{3(n-j)}+1$).
	\\
	In (2) we must have $\frac{A + 1 - c^2}{X} \in \mathbb{Z}$ which leads us to the following three cases
	\begin{enumerate}
		\item[($\alpha$)] $A + 1 - c^2 \ge X$
	\end{enumerate}
	which is impossible since $A < \sqrt{X}$.
	\begin{enumerate}
		\item[($\beta$)] $A+1-c^2=0 \Longrightarrow c = \sqrt{A+1}$
	\end{enumerate}
	In this case we will have $ab=(AX+\sqrt{A+1})(X-\sqrt{A+1})=AX^2-A \sqrt{A+1} X + \sqrt{A+1} X -(A+1)= AX(X-\sqrt{A+1}) + \sqrt{A+1}X-(A+1)$ which is clearly larger than $(X^{3/2})$, since $\sqrt{A+1}X-(A+1) > 0$ and $AX(X-\sqrt{A+1}) > X^{3/2}$ (since $ X-\sqrt{A+1} > \sqrt{X}$), hence, we have excluded this case as well.
	
	\begin{enumerate}
		\item[($\gamma$)] $c^2-A-1 \ge X$
	\end{enumerate}
	In this case we should have $A+1=c^2 \pmod X = (X-c)^2 \pmod X$, but since $(X-c) < \sqrt{X}$ then $(X-c)^2 < X$ and we can deduce that $A=(X-c)^2-1 \Longrightarrow \frac{N-1}{X^3}= ((X-c)^2-1)(X-c) + c -\frac{c^2 -1 -((X-c)^2-1)}{X} = (X-c)^3 -X +2c - \frac{c^2 -1 -X^2 +2cX -c^2 +1}{X}= (X-c)^3$, and we have excluded that case by the assumption $N \neq H^3 p^{3(n-j)}+1$. 
	\\
	Hence, we have excluded all the cases of (ii), and then we can deduce that $q=N$ and $N$ is prime.
\end{proof}

\end{document}
